%===============================================================================
% Kapitel 5: Schlussfolgerung
%===============================================================================

\chapter{Schlussfolgerung}
\label{chap:conclusion}

Die vorliegende Arbeit hat sich zum Ziel gesetzt, folgende Forschungsfragen
zu untersuchen: [Forschungsfragen anführen].

\section{Zusammenfassung der Ergebnisse}
\label{sec:conc:summary}

Die wichtigsten Erkenntnisse dieser Arbeit sind:

\begin{enumerate}
    \item \textbf{Erkenntnisse 1}: [Kurze Zusammenfassung]
    \item \textbf{Erkenntnisse 2}: [Kurze Zusammenfassung]
    \item \textbf{Erkenntnisse 3}: [Kurze Zusammenfassung]
\end{enumerate}

Diese Ergebnisse beantworten die anfängliche Forschungsfrage teilweise einfach/eindeutig.

\section{Relevanz und Bedeutung}
\label{sec:conc:relevance}

Die Erkenntnisse dieser Arbeit haben Bedeutung für:

\begin{itemize}
    \item \textbf{Wissenschaft}: Die Arbeit trägt zum Verständnis von [Konzept] bei
          und eröffnet neue Forschungsfragen.
    \item \textbf{Praxis}: Die Ergebnisse können zur Verbesserung von [Anwendungsgebiet]
          beitragen.
    \item \textbf{Gesellschaft}: Langfristig könnten die Erkenntnisse zu
              [gesellschaftlichen Veränderungen] führen.
\end{itemize}

\section{Abschließende Bemerkungen}
\label{sec:conc:remarks}

Diese Studie zeigt, dass [Kernaussage]. Obwohl Limitationen bestehen,
deuten die Befunde darauf hin, dass...

Die vorliegende Arbeit legt den Grundstein für zukünftige Forschungen,
die sich mit [verwandtem Thema] befassen sollten.

\section{Ausblick}
\label{sec:conc:outlook}

Mit den gewonnenen Erkenntnissen ergeben sich folgende nächste Schritte:

\begin{enumerate}
    \item Replikation der Studie in anderen Kontexten
    \item Untersuchung von Moderatorvariablen
    \item Entwicklung von Interventionen basierend auf den Erkenntnissen
    \item Integration der Ergebnisse in bestehende Theorien
\end{enumerate}

Insgesamt bietet diese Arbeit einen wertvollen Beitrag zum Verständnis
von [Forschungsthema] und hebt die Notwendigkeit weiterer Forschung hervor.
